\section{Method}
\label{sec:method}
Given a single reference image of a character portrait, our approach can generate a video synchronized with an input voice audio clip, preserving the natural head motion and vivid expression in harmony with the tonal variances of the provided vocal audio. By creating a seamless series of cascaded video, our model facilitates the generation of long-duration talking portrait videos with consistent identity and coherent motion, which are crucial for realistic applications.

\subsection{Preliminaries}

Our methodology employs Stable Diffusion (SD) as the foundational framework. SD is a widely-utilized text-to-image (T2I) model that evolves from the Latent Diffusion Model (LDM)\cite{ldm}. It utilizes an autoencoder Variational Autoencoder (VAE)\cite{vae} to map the original image feature distribution $x_0$ into latent space $z_0$, encoding the image as $z_0=\mathbf{E}(x_0)$ and reconstructing the latent features as $x_0=\mathbf{D}(z_0)$. This architecture offers the advantage of reducing computational costs while maintaining high visual fidelity. Based on the Denoising Diffusion Probabilistic Model (DDPM)\cite{ddpm} or the Denoising Diffusion Implicit Model (DDIM)\cite{ddim} approach, SD introduces Gaussian noise $\epsilon$ to the latent $z_0$ to produce a noisy latent $z_t$ at a specific timestep $t$. During inference, SD aims to remove the noise $\epsilon$ from the latent $z_t$ and incorporates text control to achieve the desired outcome by integrating text features. The training objective for this denoising process is expressed as:
\begin{equation}
\mathcal{L}=\mathbb{E}_{t,c, z_t, \epsilon}\left[ ||\epsilon - \epsilon_{\theta}(z_t, t, c)||^2\right]
\end{equation}
where $c$ represents the text features, which are obtained from the token prompt via the CLIP\cite{clip} ViT-L/14 text encoder. In SD, $\epsilon_{\theta}$ is realized through a modified UNet\cite{unet} model, which employs the cross-attention mechanism to fuse $c$ with the latent features.

\subsection{Network Pipelines}

\begin{figure}[tb]
  \centering
  \includegraphics[width=\textwidth]{images/pipeline.png}
  \caption{Overview of the proposed method. Our framework is mainly constituted with two stages. In the initial stage, termed Frames Encoding, the ReferenceNet is deployed to extract features from the reference image and motion frames. Subsequently, during the Diffusion Process stage, a pretrained audio encoder processes the audio embedding. The facial region mask is integrated with multi-frame noise to govern the generation of facial imagery. This is followed by the employment of the Backbone Network to facilitate the denoising operation. Within the Backbone Network, two forms of attention mechanisms are applied: Reference-Attention and Audio-Attention. These mechanisms are essential for preserving the character's identity and modulating the character's movements, respectively. Additionally, Temporal Modules are utilized to manipulate the temporal dimension, and adjust the velocity of motion.
  }
  \label{fig:pipeline}
\end{figure}

The overview of our method is shown in Figure \ref{fig:pipeline}. Our \textbf{Backbone Network} get the multi-frame noise latent input, and try to denoise them to the consecutive video frames during each time step, the Backbone Network has the similar UNet structure configuration with the SD 1.5. 1) Similar to previous work, to ensure the continuity between generated frames, the Backbone Network is embedded with \textbf{temporal modules}. 2) To maintain the ID consistency of the portrait in the generated frames, we deploy a UNet structure called \textbf{ReferenceNet} parallel to the Backbone, it input the reference image to get the features. 3) To drive the character speaking motion, \textbf{audio layers} are utilized to encode the voice features. 4) To make the motion of talking character controllable and stable, we use the \textbf{face locator} and \textbf{speed layers} to provide weak conditions.

\subsubsection{Backbone Network.} 
\label{sec:backbone}
In our work, the prompt embedding is not utilized; hence, we have adapted the cross-attention layers in the SD 1.5 UNet structure to reference-attention layers. These modified layers now take reference features from ReferenceNet as input rather than text embeddings.

\subsubsection{Audio Layers.}
\label{sec:audiolayer}
The pronunciation and tone in the voice is the main driven sign to the generated character. The features extracted from the input audio sequence by the various blocks of the pretrained wav2vec\cite{wav2vec} are concatenated to yield the audio representation embedding, $A^{(f)}$, for the $f$th frame. However, the motion might be influenced by the future/past audio segments, for example, opening mouth and inhaling before speaking. To address that, we define voice features of each generated frame by concatenating the features of nearby frames: 
$\mathbf{A}^{(f)}=\oplus\{A^{(f-m)},...A^{(f)},...A^{(f+m)}\}$, $m$  is the number of additional features from one side. To inject the voice features into the generation procedure, we add audio-attention layers performing a cross attention mechanism between the latent code and $\mathbf{A}$ after each ref-attention layers in the Backbone Network. 

% \subsubsection{ReferenceNet.} The ReferenceNet possesses a structure identical to that of the Backbone Network and serves to extract detailed features from input images. Given that both the ReferenceNet and the Backbone Network originate from the same original SD 1.5 UNet architecture, the feature maps generated by these two structures at certain layers are likely to exhibit similarities. Consequently, this facilitates the Backbone Network's integration of features extracted by the ReferenceNet. Prior research\cite{googletryon, animate_anyone} has underscored the profound influence of utilizing analogous structures in maintaining the consistency of the target object's identity. In our study, both the ReferenceNet and the Backbone Network inherit weights from the original SD UNet. The image of the target character is inputted into the ReferenceNet to extract the reference feature maps outputs from the self-attention layers. During the Backbone denoising procedure, the features of corresponding layers undergo a reference-attention layers with the extracted feature maps. As the ReferenceNet is primarily designed to handle individual images, it lacks the temporal layers found in the Backbone.

\subsubsection{ReferenceNet.} 
\label{sec:referencenet}
The ReferenceNet possesses a structure identical to that of the Backbone Network and serves to extract features from input images. Prior research\cite{googletryon, animate_anyone} has underscored the profound influence of utilizing analogous structures in maintaining the consistency of the target object's identity. In our study, both the ReferenceNet and the Backbone Network inherit weights from the original SD UNet. The reference image is inputted into the ReferenceNet to extract the reference feature maps outputs from the self-attention layers. During the Backbone denoising procedure, the features of corresponding layers undergo a reference-attention layers with the extracted feature maps.

% \subsubsection{Temporal Modules.} Most works try inserting temporal mixing layers into the pretrained text-to-image architecture, to facilitate the understanding and encoding of temporal relationships between consecutive video frames. By doing so, the enhanced model is able to maintain continuity and consistency across frames, resulting in the generation of smooth and coherent video streams. Informed by the architectural concepts of AnimateDiff, we apply self-attention temporal layers to the features within frames. Specifically, we reconfigure the input feature map $x \in \mathbb{R}^{b \times c \times f \times h \times w}$ to the shape ${(b \times h \times w) \times f \times c}$. Here, $b$ stands for the batch size, $h$ and $w$ indicate the spatial dimensions of the feature map, $f$ is the count of generated frames, and $c$ is the feature dimension. Notably, we direct the self-attention across the temporal dimension $f$, to effectively capture the dynamic content of the video. The temporal layers are inserted at each resolution stratum of the Backbone Network. 
% Most current diffusion-based video generation models are inherently limited by their design to produce a predetermined number of frames, thereby constraining the creation of extended video sequences. This limitation is particularly impactful in applications of talking head videos, where a sufficient duration is essential for the articulation of meaningful speaking. Some methodologies employ a frame from the end of the preceding clip as the initial frame of the subsequent generation, aiming to maintain a seamless transition across concatenated segments. Inspired by that, our approach incorporates the last $n$ frames, termed 'motion frames' from the previously generated clip to enhance cross-clip consistency. Specifically, these $n$ motion frames are fed into the ReferenceNet to pre-extract multi-resolution motion feature maps. During the denoising process within the Backbone Network, we merge the temporal layer inputs with the pre-extracted motion features of matching resolution along the frames dimension. This straightforward method effectively ensures coherence among various clips. For the generation of the first video clip, we initialize the motion frames as zero maps.

\subsubsection{Temporal Modules.} 
\label{sec:temporalmodule}
Informed by the architectural concepts of AnimateDiff, we apply self-attention temporal layers to the features within frames. Specifically, we reconfigure the input feature map $x \in \mathbb{R}^{b \times c \times f \times h \times w}$ to the shape ${(b \times h \times w) \times f \times c}$. Here, $b$ stands for the batch size, $h$ and $w$ indicate the spatial dimensions of the feature map, $f$ is the count of generated frames, and $c$ is the feature dimension. Notably, we direct the self-attention across the temporal dimension $f$, to effectively capture the dynamic content of the video. The temporal layers are inserted at each resolution stratum of the Backbone Network. 
Most diffusion-based video generation models are inherently limited by their design to produce a predetermined number of frames, thereby constraining the creation of extended video. This limitation is particularly impactful in applications of talking head videos, where a sufficient duration is essential for the articulation of meaningful speaking. Some methodologies employ a frame from the end of the preceding clip as the initial frame of the subsequent generation, aiming to maintain a seamless transition across concatenated segments. Inspired by that, our approach incorporates the last $n$ frames, termed 'motion frames' from the previously generated clip to enhance cross-clip consistency. Specifically, these $n$ motion frames are fed into the ReferenceNet to pre-extract multi-resolution motion feature maps. During the denoising process within the Backbone Network, we merge the temporal layer inputs with the pre-extracted motion features of matching resolution along the frames dimension. This straightforward method effectively ensures coherence among various clips. For the generation of the first video clip, we initialize the motion frames as zero maps.

It should be noted that while the Backbone Network may be iterated multiple times to denoise the noisy frames, the target image and motion frames are concatenated and input into the ReferenceNet only once. Consequently, the extracted features are reused throughout the process, ensuring that there is no substantial increase in computational time during inference.

\subsubsection{Face Locator and Speed Layers.} Temporal modules can guarantee continuity of the generated frames and seamless transitions between video clips, however, they are insufficient to ensure the consistency and stability of the generated character's motion across the clips, due to the independent generation process. Previous works use some signal to control the character motion, e.g. skeleton\cite{animate_anyone}, blendshape\cite{sadtalker}, or 3DMM\cite{sun2023vividtalk}, nevertheless, employing these control signals may be not good in creating lively facial expressions and actions due to their limited degrees of freedom, and the inadequate labeling during training stage are insufficient to capture the full range of facial dynamics. Additionally, the same control signals could result in discrepancies between different characters, failing to account for individual nuances. Enabling the generation of control signals may be a viable approach\cite{sun2023vividtalk}, yet producing lifelike motion remains a challenge. Therefore, we opt for a "weak" control signal approach. 

Specifically, as shown in Figure \ref{fig:pipeline}, we utilize a mask $\mathbf{M}=\bigcup_{i=1}^{f} M^{i}$ as the face region, which encompasses the facial bounding box (bbox) regions of the video clip. We employ the Face Locator, which consists of lightweight convolutional layers designed to encode the bounding box mask. The resulting encoded mask is then added to the noisy latent representation before being fed into the Backbone. We can use the mask to control where the character face should be generated. 

% However, creating consistent and smooth motion between clips is challenging due to variations in head motion frequency during separate generation processes. To address this issue, we incorporate the target head motion speed into the generation. More precisely, we consider the head rotation velocity $w^f$ in frame $f$ and divide the range of speeds into $d$ discrete speed buckets, each representing a different velocity level. Each bucket has a central value $c^d$ and a radius $r^d$. We retarget $w^f$ to a vector $S=\{s^d\}\in\mathbb{R}^d$, where $s^d=\tanh((w^f-c^d)/r^d*3)$. Similar to the method used in the audio layers, the head rotation speed embedding for each frame is given by $S^{f}=\oplus\{S^{(f-m)},\ldots,S^{(f)},\ldots,S^{(f+m)}\}$, and $S^{f}\in\mathbb{R}^{b\times f \times c^{speed}}$ is subsequently processed by an MLP to extract speed features. Within the temporal layers, we repeat $S^{f}$ to the shape ${(b \times h \times w) \times f \times c^{speed}}$ and implement a cross-attention mechanism that operates between the speed features and the reshaped feature map across the temporal dimension $f$. By doing so and specifying a target speed, we can synchronize the rotation speed and frequency of the generated character's head across different clips. Combined with the facial position control provided by the Face Locator, the resulting output can be both stable and controllable.

However, creating consistent and smooth motion between clips is challenging due to variations in head motion frequency during separate generation processes. To address this issue, we incorporate the target head motion speed into the generation. More precisely, we consider the head rotation velocity $w^f$ in frame $f$ and divide the range of speeds into $d$ discrete speed buckets, each representing a different velocity level. Each bucket has a central value $c_i\in \{c_1,...,c_d\}$ and a radius $r_i \in \{r_1, ..., r_d\}$. We retarget $w^f$ to a vector $\mathbf{s}\in \mathbb{R}^d$, where the $i$th value notated by $s_i = \tanh((w^f-c_i)/r_i*3)$.
% $S=\{s^d\}\in\mathbb{R}^d$, where $s^d=\tanh((w^f-c^d)/r^d*3)$. 
Similar to the method used in the audio layers, the head rotation speed embedding for each frame is given by $\mathbf{S}^{f}=\oplus\{\mathbf{s}^{(f-m)},\ldots,\mathbf{s}^{(f)},\ldots,\mathbf{s}^{(f+m)}\}$. The speed embedding of each clip denoted by $\mathbf{S}\in\mathbb{R}^{b\times f \times (2m+1)d}$ is subsequently processed by an MLP into a speed feature map $\mathbf{F}\in\mathbb{R}^{b\times f\times l}$. Within the temporal layers, we repeat $\mathbf{F}$ to the shape ${(b \times h \times w) \times f \times l}$ and implement a cross-attention mechanism that operates between the speed features and the reshaped feature map across the temporal dimension $f$. By doing so and specifying a target speed, we can synchronize the rotation speed and frequency of the generated character's head across different clips. Combined with the facial position control provided by the Face Locator, the resulting output can be both stable and controllable.


It should also be noted that the specified face region and assigned speed does not constitute strong control conditions. In the context of face locator, since the $\mathbf{M}$ is the union area of the entire video clip, indicating a sizeable region within which facial movement is permissible, thereby ensuring that the head is not restricted to a static posture. With regard to the speed layers, the difficulty in accurately estimating human head rotation speed for dataset labeling means that the predicted speed sequence is inherently noisy. Consequently, the generated head motion can only approximate the designated speed level. This limitation motivates the design of our speed buckets framework. 


\subsection{Training Strategies}

The training process is structured into three stages. The first stage is the image pretraining, where the Backbone Network, the ReferenceNet, and the Face Locator are token into training, in this stage, the Backbone takes a single frame as input, while ReferenceNet handles a distinct, randomly chosen frame from the same video clip. Both the Backbone and the ReferenceNet initialize weights from the original SD. In the second stage, we introduce the video training, where the temporal modules and the audio layers are incorporated, $n + f$ contiguous frames are sampled from the video clip, with the started $n$ frames are motion frames. The temporal modules initialize weights from AnimateDiff\cite{animatediff}. In the last stage, the speed layers are integrated, we only train the temporal modules and the speed layers in this stage. This strategic decision deliberately omits the audio layers from the training process. Because the speaking character's expression, mouth motion, and the frequency of the head movement, are predominantly influenced by the audio. Consequently, there appears to be a correlation between these elements, the model might be prompted to drive the character's motion based on the speed signal rather than the audio. Our experimental results suggest that simultaneous training of both the speed and audio layers undermines the driven ability of the audio on the character's motions.